\subsection{Problem Statement}
\label{sec:problem-statement}

We approach intrusion detection as a \emph{supervised, flow-level, multiclass classification} problem solved in near real time. Let $\mathcal{F}$ denote the population of bidirectional network flows produced by the \textsc{CICFlowMeter} extractor on a monitored network. Each flow $\mathbf{x} \in \mathcal{F}$ is summarised by a feature vector $\mathbf{x} \in \mathbb{R}^{d}$ with $d = 78$ in the stock CIC-IDS2017 schema (extended in Section~\ref{sec:feat-eng}). The label space is $\mathcal{Y} = \{0, 1, \ldots, C{-}1\}$, where class $0$ denotes \texttt{BENIGN} traffic and the remaining $C{-}1$ classes correspond to the aggregate attack categories enumerated in Table~\ref{tab:class-dist}.

The IDS must learn a classifier $f_\theta : \mathbb{R}^{d} \to \mathcal{Y}$ from a labelled training corpus $\mathcal{D}_{\text{train}} \subset \mathcal{F} \times \mathcal{Y}$ that minimises the expected risk on the underlying flow distribution. Two operating regimes coexist:

\begin{description}[leftmargin=1.2em,itemsep=2pt,topsep=2pt]
    \item[Offline training.] $f_\theta$ is fitted on a static, fully labelled snapshot. The training step has no real-time constraint and is allowed to spend tens of minutes to a few hours of compute, including stratified cross-validation and hyperparameter search.
    \item[Online inference.] At deployment, $f_\theta$ receives a continuous stream of flows $\mathbf{x}_1, \mathbf{x}_2, \ldots$ produced by the live flow extractor. For each incoming $\mathbf{x}_t$ the system must emit a prediction $\hat{y}_t = f_\theta(\mathbf{x}_t)$ within a strict per-flow latency budget so that downstream alert and blocking actions remain timely.
\end{description}

\subsubsection*{Operational Constraints.}
The proposed IDS targets the following requirements, derived from typical inline deployment scenarios on enterprise edge gateways:

\begin{description}[leftmargin=1.2em,itemsep=2pt,topsep=2pt]
    \item[Detection quality.] Macro-$F_1$ over all $C$ classes shall be at least $0.95$ on the held-out test set, ensuring that minority attack categories such as \texttt{Bot} or \texttt{Web Attack} are not silently traded for benign accuracy.
    \item[False-positive rate (FPR).] FPR on benign flows shall not exceed $1\%$. Above this threshold the alert volume overwhelms human analysts and produces alert fatigue, which is the dominant operational failure mode of real-world IDS deployments.
    \item[Per-flow latency.] End-to-end inference latency---measured from the moment a complete flow record is emitted by the flow extractor to the moment the prediction is available to downstream consumers---shall not exceed $10\,\text{ms}$ at the $99^{\text{th}}$ percentile on commodity hardware (8-core x86 CPU, no GPU).
    \item[Throughput.] Sustained throughput shall reach at least $10^4$ flows per second on a single inference worker, which is sufficient for a 10\,Gbps link with average flow sizes representative of CIC-IDS2017.
    \item[Resource footprint.] The serialised model, the scaler/encoder artefacts, and the runtime feature buffers shall together occupy less than $512\,\text{MB}$ of resident memory, allowing co-deployment with the flow extractor on the same gateway host.
\end{description}

\subsubsection*{Assumptions and Non-Goals.}
We assume that the upstream flow extractor produces complete, well-formed feature records with bounded skew between training and serving (\emph{no} concept drift handling is performed in the present version; see Section~\ref{sec:bg-network} for related discussion). We further assume that the training set is representative of the operational traffic mix; significant distribution shift would invalidate the offline-fitted model and motivate the continual-learning extensions discussed in our concluding section. Active adversarial evasion---where attackers deliberately craft flows to be misclassified---is acknowledged but treated as out of scope for this report.
